\documentclass[10pt,a4paper]{amsart}
\usepackage[margin=19mm]{geometry}
\usepackage{amsmath,amssymb,mathtools}
\usepackage[scr=rsfs,cal=euler]{mathalfa}
\usepackage{xfrac}
\usepackage[unicode,hidelinks,bookmarks=false]{hyperref}
\newcommand{\ie}{i.e.\ }
\newcommand{\cf}{cf.\ }
\newcommand{\resp}{respectively\ }
\newcommand{\Subsection}{\subsection}
\providecommand{\nobreakdash}{\nobreak\hbox{-}}
\setlength{\emergencystretch}{2em}
\multlinegap0pt
\usepackage{bm}
\usepackage{bbm}
\usepackage{dsfont}
\usepackage{xspace}
\usepackage{cancel}
\usepackage{mathtools}
\usepackage{amsthm}
\makeatletter
\@ifundefined{theorem}{\newtheorem{theorem}{Theorem}}{}
\@ifundefined{coro}{\newtheorem{coro}[theorem]{Corollary}}{}
\@ifundefined{lemma}{\newtheorem{lemma}[theorem]{Lemma}}{}
\makeatother
\newcommand{\dam}{\beta} 
\newcommand{\e}{\varepsilon} 
\title[Approach to equilibrium for a particle interacting with a ha]{Approach to equilibrium for a particle interacting with a harmonic thermal bath}
\author{Federico Bonetto, Alberto Mario Maiocchi}
\date{Source: arXiv:2510.20003v2 (CC BY 4.0)}
\begin{document}
\maketitle

\noindent\emph{Source and licence.} Approach to equilibrium for a particle interacting with a harmonic thermal bath. Version arXiv:2510.20003v2 (CC BY 4.0), distributed under the Creative Commons Attribution 4.0 International licence, as recorded in the arXiv licence metadata for this exact version and shown on the abstract page https://arxiv.org/abs/2510.20003v2. Full rights notice: licenses/arxiv-2510.20003-CC-BY-4.0.txt.\par\smallskip

\noindent\textit{Editorial scope.} Counted unit: the Lemma whose source label is \texttt{lem:duhamel} in the article (source0.tex lines 2474--2489, source label \texttt{lem:duhamel}), delivered together with the article's own complete original proof of it (source0.tex lines 2491--2547, one \texttt{proof} environment of 57 lines). No mathematics is written, repaired or re-derived: statement and proof are the article's own text line for line. Required context retained uncounted: lem:stima (lines 2413--2421); cor:stimasing (lines 2448--2456). Every definition, display and label of those lines is retained unchanged. Omitted from this package and not redistributed: 1--2412, 2422--2447, 2457--2473, 2490--2490, 2548--2774. Nothing inside a retained excerpt is omitted or altered: every retained body is delivered exactly as the article's lines are written, so no omission receipt of any kind accompanies this package. Citations: the retained excerpts carry no citation command at all, so no bibliography entry of the article is transcribed and no reference is redistributed. Reference adaptation: none was needed and none was made; every reference inside the delivered unit resolves inside it, so no unresolved reference or citation is produced. Printed slips kept literal: no printed word, symbol, hypothesis or number of the counted statement or of its proof was altered. Layout and class: the article's own class is \texttt{article} with packages including latexsym, amsfonts, amsmath, amsthm, mathrsfs, dsfont, bbold, babel, caption, epsfig, float, subfigure, psfrag, graphicx; the wrapper is the lane's pinned \texttt{amsart} editorial wrapper at 10pt on A4, which restates the theorem-like environments the retained text needs and adds the article's own macro definitions that the retained text needs (\texttt{\textbackslash dam}, \texttt{\textbackslash e}), with extra wrapper packages (bm, bbm, dsfont, xspace, cancel, mathtools). The delivered numbering is the unit's own. Assets: no retained span contains a figure environment, a graphics-inclusion command, a PDF-page inclusion, a table or a TikZ picture; no image file is copied into this package, the package declares no asset dependency and no image file is redistributed with it. Licence: version arXiv:2510.20003v2 of this article is distributed under the Creative Commons Attribution 4.0 International licence, as recorded in the arXiv licence metadata for this exact version and shown on the abstract page https://arxiv.org/abs/2510.20003v2. A full rights notice is delivered at licenses/arxiv-2510.20003-CC-BY-4.0.txt. Characters: every retained excerpt is pure ASCII, so the delivered engine, XeLaTeX with its default Latin Modern text font, reports no missing character.
\begin{lemma}\label{lem:stima} 
Let $h(\xi)$ be analytic in $\mathcal R=\{\xi; |\Re(\xi)|< 1\}$ and for $d>0$, let $\mathcal R_d=\{\xi\in \mathcal R;
|\Im(\xi)|\leq d\}$. Then, for $\e>-1$ and $t\ge 0$ we have
\begin{equation}\label{eq:fint}
\left |\int_{-1}^1 (1-\xi^2)^{\e}h(\xi)e^{i \xi t}d\xi\right|\leq
\frac{K_{d,\e}\sup_{\xi\in \mathcal R_d}|h(\xi)|}{\left(1+dt\right)^{1+\e}}\ .
\end{equation}

\end{lemma}

\begin{coro}\label{cor:stimasing}
 Let $h(\xi)$ be analytic in $\mathcal R$. Then for $\zeta\in \mathcal R$ with $0<\Im(\zeta)\leq 1$ and $\e>-1$ we have
\begin{equation}\label{eq:fintsing}
\left |\int_{-1}^1 \frac{(1-\xi^2)^{\e}h(\xi)}{\xi-\zeta}e^{i
\xi t}d\xi-2\pi i(1-\zeta^2)^\e h(\zeta)e^{i\zeta t}\right|\leq
\frac{K'_{\e}\sup_{\xi\in
\mathcal R_{2/(1-|\Re(\zeta)|)}}|h(\xi)|}{(1-|\Re{\zeta}|)^{2+\epsilon}(1+t)^{1+\e}}\ .
\end{equation}
\end{coro}

\begin{lemma}\label{lem:duhamel}
Let $0<a<b$, let $\dam$, $\Xi$ be such that $\Xi\in(a,b)$, while $0<\dam\leq \min\{\Xi-a,b-\Xi\}/2$ and let $g(\omega)$ be a function real for
$\omega\in [a,b]$ and analytic in $\mathcal R_{a,b}= \{\xi: a\le
\Re(\xi)\le b, |\Im(\xi)|\le \max\{1,\dam\}\}$. Then we have,
\begin{equation}\label{eq:duha}
\begin{aligned}
 \int_{a}^{b}d\omega\, &\frac{g(\omega)}{\sqrt{(\omega-a)(b-\omega)}} 
 \int_0^t\int_0^s\cos(\omega(t-s-\tau+\sigma))
 e^{-\dam \tau}\sin(\Xi\tau)e^{-\dam\sigma}\sin(\Xi\sigma)d\sigma d\tau=\\
&\frac{\pi}{4\dam}\frac{g(\Xi)}{\sqrt{(\Xi-a)(b-\Xi)}}  
\left(e^{-\dam|t-s|}-e^{-\dam(t+s)}\right)\cos(\Xi(t-s))+
\frac{K(t,s)}{\left((\Xi-a)(b-\Xi)\right)^{2+\e}}\ ,
\end{aligned}
\end{equation}
where $K(t,s)$ is bounded uniformly in $\beta$, $s$ and $t$.
\end{lemma}

\begin{proof} Observe that
\begin{equation}\label{eq:contr1}
\begin{aligned}
\int_0^td\tau\int_0^sd\sigma&
\cos(\omega(t-s-\tau+\sigma))
\sin(\Xi\tau)e^{-\dam \tau}\sin(\Xi\sigma)e^{-\dam \sigma}=\\
&-\frac 18\sum_{g_1,g_2,g_3=\pm}g_2g_3\int_0^td\tau\int_0^sd\sigma
g_2g_3e^{g_1i\omega(t-s-\tau+\sigma)}
e^{(g_2i\Xi-\dam)\tau}
e^{(g_3i\Xi-\dam)\sigma}=\\
&-\frac 18\sum_{g_1,g_2,g_3=\pm}g_2g_3e^{g_1i\omega(t-s)}
\int_0^t e^{[i(-g_1\omega+g_2\Xi)-\dam]\tau}d\tau
\int_0^s e^{[i(g_1\omega+g_3\Xi)-\dam]\sigma}d\sigma=\\
&-\frac 18\sum_{g_1,g_2,g_3=\pm}g_2g_3e^{g_1i\omega(t-s)}
\frac{1-e^{[i(-g_1\omega+g_2\Xi)-\dam]t}}
{i(-g_1\omega+g_2\Xi)-\dam}\;
\frac{1-e^{[i(g_1\omega+g_3\Xi)-\dam]s}}
{i(g_1\omega+g_3\Xi)-\dam}=\\
&-\frac 18\sum_{g_1,g_2,g_3=\pm}g_2g_3
\frac{e^{ig_1\omega t}-e^{[ig_2\Xi-\dam]t}}
{\omega-g_1g_2\Xi-ig_1\dam}\;
\frac{e^{-ig_1\omega s}-e^{[ig_3\Xi-\dam]s}}
{\omega+g_1g_3\Xi+ig_1\dam}\ .
\end{aligned}
\end{equation}
Notice that the terms for  $g_2=-g_1$ and $g_3=g_1$ are bounded
uniformly in $\omega\in[a,b]$ and $\dam$,$s$ and $t$. If $g_2=g_3$ the
corresponding term have one pole
close to $[-1,1]$ but the residue is 
bounded
by a constant independent of $\dam$. Using Corollary \ref{cor:stimasing} we obtain that also contribution of this
term to \eqref{eq:duha} can be bounded uniformly in $\dam$, $s$ and $t$. 

We are thus left with the contribution for $g_2=g_1=-g_3$, 
\begin{equation}\label{eq:contazzo}
\begin{aligned}
&\frac 18\sum_{g_1=\pm}\int_{a}^{b}d\omega\, \frac{g(\omega)}{\sqrt{(\omega-a)(b-\omega)}}
\frac{\left(e^{ig_1\omega t}-e^{[ig_1\Xi-\dam]t}\right)\left(e^{-ig_1\omega 
s}-e^{[-ig_1\Xi-\dam]s}\right)}{(\omega-\Xi)^2+\beta^2}=\\
&\frac 18\int_{a}^{b}d\omega\, \frac{g(\omega)}{((\omega-\Xi)^2+\dam^2) \sqrt{(\omega-a)(b-\omega)}}\\
&\qquad\qquad\qquad\qquad\sum_{g_1=\pm}\left(e^{ig_1\omega(t-s)}-e^{(ig_1\Xi-\dam)t-ig_1\omega s}-e^{ig_1\omega t - 
(ig_1\Xi+\dam)s}+
e^{ig_1\Xi(t-s)-\dam(t+s)}\right)=\\
&\frac {\pi}{4\dam}\frac{g(\Xi)}{\sqrt{(\Xi-a)(b-\Xi)}}e^{-\dam|t-s|}\cos(\Xi(t-s))+\\
&\qquad\qquad
-\frac
\pi{4\dam}\frac{g(\Xi)}{\sqrt{(\Xi-a)(b-\Xi)}}
e^{-\dam(t+s)}\cos(\Xi(t-s))+\frac{K(t,s)}{\left((\Xi-a)(b-\Xi)\right)^{2+\e}}\ , 
\end{aligned}
\end{equation}
where $K(t,s)$ is uniformly bounded in $\dam$, $t$, and $s$ and vanishes when $t$ or $s$ go to infinity. In 
\eqref{eq:contazzo} we have expanded the product to be able to apply Lemma \ref{lem:stima} and Corollary 
\ref{cor:stimasing}. Indeed, in each term, whether to move the $\omega$ integration path for $\Im(\omega)$ positive or 
negative depends on the sign of the factor multiplying $i\omega$ in
the exponent.

\end{proof}

\end{document}
